\section{Normalization and successive attaching classes}
\label{sec:normalization}

We now use the vanishing in Equation~\eqref{eq:k-test} to construct
the cycles $b_1,\ldots,b_n$.  The vanishing is obtained only after
tensoring with a Koszul algebra.  Accordingly, we first pass to a
localization in which those tests detect equivalences, and then show
that the actual completed targets already belong to that localization.
This produces algebra maps, together with the boundary choices needed
to attach the next cell.  No identification of the cobordism Hurewicz
images is used in this construction.

\subsection{The inductive data}

Retain the lattices $\Lambda$, $\Omega$, and
$I=\Lambda\oplus\Lambda\oplus\Z$, the homomorphism $F:I\to\Omega$,
and the functions $\nu$ and $d$ from the preceding sections.  Set
\[
 i_j=(w_j,0,-1)\in I \qquad (0\leq j<n),
 \qquad \tau=(0,0,1).
\]
In particular, $F(i_j)=w_{j+1}$.  Let
$M=(\MU_{(p)})_p\in\cC$ and put
\[
 L_0=\widehat B(\one),\qquad L_0^M=\widehat B(M).
\]
Here the hats denote derived completion with respect to $t$, as in
Lemma~\ref{layer:completion}.  Starting with
$b_0=\varpi u_0$, we have the algebra
$h_\Omega(Y_1)=A[e_0]$, with $\partial e_0=b_0$.

At stage $0\leq a\leq n$, the inductive data will consist of the
strict cycle models through $b_a$ and $Y_{a+1}$, and a tower of
$I$-graded commutative algebras
\[
 L_a=L_0[e'_0,\ldots,e'_{a-1}],\qquad
 \wt(e'_j)=i_j,\quad |e'_j|_{\mathrm{chain}}=1,
 \quad \partial e'_j=g_j.
\]
This notation means successive algebra attachments with their
specified nulls.  In particular, it does not require the maps used to
construct the $g_j$ to have been prescribed as strict dg maps.  Define
the parallel tower by base change:
\[
 L_a^M=L_0^M\otimes_{L_0}L_a.
\]
Its attached classes and nulls are the images of those in $L_a$.
All tensor products of algebras in this section are derived.

The invariant is a pair of equivalences of $I$-graded commutative
algebras
\begin{equation}
\label{eq:target-layer}
 \begin{aligned}
 L_a//t&\simeq F^*h_\Omega(Y_{a+1}),\\
 L_a^M//t&\simeq F^*h_\Omega(M\otimes Y_{a+1}).
 \end{aligned}
\end{equation}
These equivalences are compatible with the sphere-to-$M$ unit, with
the maps from the initial layers, and with the algebra maps in the
$Y$-tower.  In the $\delta$ direction, their scalar action is ordinary
evaluation followed by the indicated algebra unit.  The case $a=0$
is supplied by Corollary~\ref{layer:first-koszul} and the equality of the
finite reductions of $B$ and $\widehat B$.

Each null attachment has a two-term underlying module by
Lemma~\ref{found:free-null}.  Thus $L_a$ is perfect over $L_0$, and
$L_a^M$ is perfect over $L_0^M$.  Lemma~\ref{layer:completion} shows
that both towers remain derived $t$-complete.

\subsection{Localization by the Koszul algebra}

Fix a stage $a<n$ for which these data have been constructed.
At a weight $\lambda\in\Lambda$, the power on the universal
coefficients takes values in
\[
 \bigl(B(R_p^a)[t^{-1}]\bigr)_{(\lambda,0,0)}.
\]
We seek an algebra map from $R^a$ into the components
\[
 \bigl\{(L_a)_{(\lambda,0,-\nu(\lambda))}\bigr\}_{\lambda\in\Lambda}.
\]
These components form a graded algebra because their index map is
additive. After inverting $t$ and multiplying its weight-$\lambda$
component by $t^{\nu(\lambda)}$, the new map must recover the power
followed by evaluation $R_p^a\to S_{(p)}$ and the map
$B(\one)\to L_a$. At weight $w_a$, the new target index is $i_a$.
Extending this coefficient map across the older Koszul nulls will
then let us apply it to $b_a^u$ and reduce modulo $t$ to construct
$b_{a+1}$. Thus the normalization must preserve algebra maps and
their specified nulls, rather than only divide individual classes
after inversion.

For $m>-\nu(\lambda)$, Equation~\eqref{eq:k-test} makes multiplication
by $t$ into the row $(\lambda,*,m)$ an equivalence after tensoring
with $K^a$. The following completion construction makes those
comparisons invertible; we will then show that the target $L_a$ is
unchanged by it.

For a commutative graded algebra $R$ and a commutative
$R$-algebra $K$, call an $R$-module $V$ \emph{$K$-acyclic} if
$K\otimes_R V=0$.  An $R$-module $U$ is \emph{$K$-local} if
$\map_R(V,U)=0$ for every $K$-acyclic $V$.  These definitions refer
to graded modules and their graded module category; any additional
indices can be kept as parameters.

\begin{lemma}[Perfect algebra completion]
\label{norm:amitsur}
Let $R$ be a commutative graded rational algebra, and let $K$ be a
commutative $R$-algebra whose underlying $R$-module is perfect.  Define
\[
 \mathcal T(U)=
 \Tot\bigl(K^{\otimes_R(\bullet+1)}\otimes_R U\bigr),
 \qquad c_U:U\longrightarrow\mathcal T(U),
\]
using the Amitsur coaugmentation.  Then $\mathcal T$ is exact,
$\mathcal T(U)$ is $K$-local, and $c_U$ becomes an equivalence after
tensoring with $K$.  Consequently $c_U$ is an equivalence whenever
$U$ is $K$-local, and $\mathcal T$ sends every $K$-acyclic module to
zero.  On commutative $R$-algebras the construction and its
coaugmentation are maps of commutative algebras.
\end{lemma}

\begin{proof}
The cosimplicial structure inserts the unit of $K$ in the coface maps
and multiplies adjacent $K$ factors in the codegeneracy maps.
Tensoring with $K$ commutes with totalization because a perfect module
is dualizable.  After that tensor, the augmented cosimplicial object
is split: multiplication into the extra, distinguished $K$ factor
gives the extra codegeneracy.  It follows that
\[
 K\otimes_R U\ \xrightarrow{\simeq}\
 K\otimes_R\mathcal T(U).
\]

Every nonaugmented term is a $K$-module.  If $W$ is a $K$-module and
$V$ is $K$-acyclic, extension and restriction of scalars give
\[
 \map_R(V,W)\simeq\map_K(K\otimes_R V,W)=0.
\]
Thus those terms are $K$-local.  Local objects are closed under
limits, so their totalization is local as well.  The fiber of $c_U$
is $K$-acyclic by the preceding split calculation.  If $U$ is local,
that fiber is also local, and hence is zero: its endomorphism spectrum
vanishes by applying the definition of locality to the fiber itself.
This proves the assertion about local inputs.  If $U$ is acyclic,
each term in its totalization is zero.

Tensoring over $R$ is exact in the stable module category, and limits
preserve finite limits.  Therefore $\mathcal T$ is exact.  Finally,
for algebra input the cosimplicial terms are commutative algebras and
all structure maps are algebra maps.  Their limit has the asserted
algebra structure and coaugmentation.  This is the usual
dualizable-algebra completion argument; compare
\cite[Example~2.18 and Proposition~2.21]{MNN2017}.
\end{proof}

Use the universal algebras and primewise models of
Proposition~\ref{universal:primewise-models}, and abbreviate
\[
 R=R^a,\qquad K=K^a,\qquad B=B(R_p^a).
\]
The small-weight comparison identifies $R$ with $h_\Lambda(R_p^a)$.
It acts on $B$ in the $\delta$ direction.  Evaluation of the
primewise universal algebra gives
\[
 v:B\longrightarrow B(\one)\longrightarrow L_0
       \longrightarrow L_a.
\]
The $R$-module $K$ is perfect by Lemma~\ref{universal:filtered-cells}.
We apply Lemma~\ref{norm:amitsur} with $R$ and $K$ supported at
indices $(0,\delta,0)$ of $I$.  Write the resulting functor as
$\mathcal T_a$.  Its module statements apply separately to each
row with fixed $(\lambda,m)$, while its algebra construction retains
all the $I$-graded products.

\begin{lemma}[Locality of the completed target]
\label{norm:target-locality}
Every $(\lambda,m)$ row of $L_a$ is $K^a$-local for the evaluated
$R^a$-action.  Thus the coaugmentation
$c_{L_a}:L_a\to\mathcal T_a L_a$ is an equivalence of
$I$-graded algebras.
\end{lemma}

\begin{proof}
Under Equation~\eqref{eq:target-layer}, restriction of the right-hand
algebra to $(0,\delta,0)$ is
\[
 h_\Lambda(Y_{a+1})
   =\bigl(A[e_0,\ldots,e_a]\bigr)|_\Lambda.
\]
There is an algebra map from $K^a$ to this restriction: send the
universal coefficients to the chosen coefficients in $A|_\Lambda$,
and send each $e_j$ to the actual null cell $e_j$.  Its differential
identities are exactly the universal-cycle identities after
evaluation.  The scalar compatibility in
Equation~\eqref{eq:layer} and the chosen path matching primewise and
actual coefficient evaluations identify its restriction to $R^a$
with the $R^a$-action on $L_a//t$.  Transporting along
Equation~\eqref{eq:target-layer} therefore makes every row of that
layer a $K^a$-module, and hence a local module.

The cofiber of a composite of $s$ multiplications by $t$ is a finite
extension of shifts of the cofiber of one multiplication by $t$.
Applied row by row, this shows that $L_a/t^s$ is local for every
$s\geq1$.  Here $/$ denotes the module cofiber; its underlying
module agrees with that of the corresponding algebra quotient.
Local modules are closed under finite extensions and inverse limits.
Since $L_a$ is derived $t$-complete, it is the inverse limit of these
local finite quotients and is itself local.  The conclusion follows
from Lemma~\ref{norm:amitsur}.  This proof uses the already attached
cells, and makes no assumption about any Hurewicz image.
\end{proof}

\subsection{The normalized algebra map}

For an $I$-graded algebra $U$ equipped with the parameter $t$, write
\[
 U^\#_\lambda=U_{(\lambda,0,-\nu(\lambda))},
 \qquad
 U^\flat_\lambda=U_{(\lambda,0,0)}.
\]
Both restrictions are $\Lambda$-graded algebras because their index
maps are additive.  There is a natural algebra map
\begin{equation}
\label{norm:regrading-map}
 s_U:U^\#\longrightarrow\bigl(U[t^{-1}]\bigr)^\flat,
 \qquad (s_U)_\lambda=t^{\nu(\lambda)}.
\end{equation}
In this formula one first maps to the localization and then
multiplies by the indicated Laurent power.  The powers, including
negative powers, form a coherent multiplicative family: pull back
the homogeneous units of the rational Laurent polynomial algebra
along $\lambda\mapsto\nu(\lambda)\tau$ and use the specified
algebra map into $U[t^{-1}]$.  Additivity of $\nu$ and
commutativity of multiplication prove that $s_U$ is an algebra map.

\begin{lemma}[The boundary row]
\label{norm:normalization}
For $B=B(R_p^a)$, the map
\[
 s_{\mathcal T_a B}:
 (\mathcal T_a B)^\#\longrightarrow
 \bigl((\mathcal T_a B)[t^{-1}]\bigr)^\flat
\]
is an equivalence of $\Lambda$-graded rational algebras.
\end{lemma}

\begin{proof}
By Equation~\eqref{eq:k-test}, the cofiber of multiplication by $t$
into a row $(\lambda,*,m)$ is $K$-acyclic whenever
$m> -\nu(\lambda)$.  Exactness of $\mathcal T_a$ sends this
cofiber to zero.  Multiplication by the coaugmented $t$ on
$\mathcal T_a B$ is the image under $\mathcal T_a$ of multiplication
by $t$ on $B$: this holds termwise in its cosimplicial construction
and therefore after totalization.  Consequently the transitions
\[
 (\mathcal T_a B)_{(\lambda,\delta,m-1)}
   \xrightarrow{t}
 (\mathcal T_a B)_{(\lambda,\delta,m)}
\]
are equivalences whenever $m> -\nu(\lambda)$.

The telescope for localization is thus constant up to equivalence
starting at the boundary row $m=-\nu(\lambda)$.  In particular,
that row maps equivalently to the same row after localization.
Multiplication by $t^{\nu(\lambda)}$ identifies the latter row
with the localized row at $m=0$.  This proves the assertion in every
weight.  When $\nu(\lambda)<0$, the last multiplication is simply
a negative Laurent power; the telescope still reaches the same
invertible tail.  Multiplicativity follows from
Equation~\eqref{norm:regrading-map}.
\end{proof}

We can now define the normalized operation.  The power map of
Proposition~\ref{power:operation}, the coaugmentation, and the two
equivalences just proved give
\begin{equation}
\label{norm:normalized-operation}
\begin{aligned}
 P^a:R^a\simeq h_\Lambda(R_p^a)
 &\xrightarrow{P_{R_p^a}}\bigl(B[t^{-1}]\bigr)^\flat
 \xrightarrow{(c_B[t^{-1}])^\flat}
       \bigl((\mathcal T_a B)[t^{-1}]\bigr)^\flat\\
 &\xrightarrow{s_{\mathcal T_a B}^{-1}}
       (\mathcal T_a B)^\#
 \xrightarrow{(\mathcal T_a v)^\#}
       (\mathcal T_a L_a)^\#
 \xrightarrow{(c_{L_a}^\#)^{-1}} L_a^\#.
\end{aligned}
\end{equation}
This is a map of graded algebras over rational spectra.  The power
map can change the action of the varying scalar ring $D$, and
$P^a$ is not asserted to be $D$-linear.

\begin{lemma}[Recovery and naturality]
\label{norm:naturality}
The normalized map recovers the evaluated power after inversion:
\begin{equation}
\label{norm:raw-recovery}
 s_{L_a}P^a\simeq (v[t^{-1}])^\flat P_{R_p^a}.
\end{equation}
For $a>0$, its restriction to $R^{a-1}$ agrees, as a rational
algebra map, with $P^{a-1}$ followed by the map $L_{a-1}^\#\to
L_a^\#$.
\end{lemma}

\begin{proof}
The first assertion says that the following square commutes:
\[
\begin{array}{ccc}
 R^a&\xrightarrow{\ P^a\ }&L_a^\#\\[2pt]
 {\scriptstyle P_{R_p^a}}\downarrow&&
       \downarrow{\scriptstyle s_{L_a}}\\[2pt]
 (B[t^{-1}])^\flat&\xrightarrow{\ (v[t^{-1}])^\flat\ }&
       (L_a[t^{-1}])^\flat .
\end{array}
\]
To verify it, naturality of $s$ for $c_{L_a}$ and $\mathcal T_a v$,
followed by naturality of the coaugmentation for $v$, gives
\[
\begin{aligned}
 &s_{L_a}(c_{L_a}^\#)^{-1}(\mathcal T_a v)^\#
        s_{\mathcal T_a B}^{-1}(c_B[t^{-1}])^\flat\\
 &\qquad=\bigl((c_{L_a}[t^{-1}])^\flat\bigr)^{-1}
       \bigl((\mathcal T_a v)[t^{-1}]\bigr)^\flat
       (c_B[t^{-1}])^\flat
   \simeq (v[t^{-1}])^\flat.
\end{aligned}
\]
The inverse in this display is the localization of an already
invertible coaugmentation on $L_a$.

For the second assertion, the compatible universal models and
evaluations give a map from the preceding diagram
$K^{a-1}\leftarrow R^{a-1}\to B(R_p^{a-1})\to L_{a-1}$
to the current diagram $K^a\leftarrow R^a\to B(R_p^a)\to L_a$.
In cosimplicial degree $r$, it induces an algebra map
\[
 (K^{a-1})^{\otimes_{R^{a-1}}(r+1)}
       \otimes_{R^{a-1}}B(R_p^{a-1})
 \longrightarrow
 (K^a)^{\otimes_{R^a}(r+1)}\otimes_{R^a}B(R_p^a).
\]
The same construction applies with $L_{a-1}$ and $L_a$ in the last
positions.  These maps commute with the units and multiplications
defining the cosimplicial structure, so they induce maps of
coaugmented totalizations.  They also preserve $t$.

The maps $s$ are natural for these maps of totalizations and for their
localizations.  Both normalization maps being inverted in
Equation~\eqref{norm:normalized-operation} are equivalences, so
their inverses respect these naturality squares.  The same is true
of the two target coaugmentations.  Finally, $P_{R_p^a}$ is natural
for the map of actual primewise algebras, and the chosen small-weight
comparisons and evaluation paths extend those of the preceding
stage.  Comparing all the arrows in
Equation~\eqref{norm:normalized-operation} therefore proves the
claimed homotopy of algebra maps.  This comparison uses actual maps
between the two Amitsur diagrams.  It does not interchange
totalization with $t$-inversion or identify a totalization with its
extension of scalars.
\end{proof}

\subsection{Extending across nulls and preserving the layer}

The map $P^a$ acts on universal coefficients.  To apply it to the
universal cycle $b_a^u$, we must also specify images of its older
Koszul nulls.  We construct
\[
 f^a:K_-^a\longrightarrow L_a^\#
\]
extending $P^a$.  For $a=0$, set $f^0=P^0$, since $K_-^0=R^0$.
For $a>0$, there is a pushout description
\[
 K_-^a\simeq R^a\otimes_{R^{a-1}}K^{a-1}.
\]
By induction, $f^{a-1}$ gives the map on $K_-^{a-1}$, followed by
$L_{a-1}^\#\to L_a^\#$.  The attached cell $e'_{a-1}$ in $L_a$
is a specified null of
$g_{a-1}=f^{a-1}(b_{a-1}^u)$.  Its index
$(w_{a-1},0,-1)$ is exactly the index selected by the $\#$ restriction
at $w_{a-1}$.  The universal property of the null attachment extends
the map to $K^{a-1}$.  Its restriction to $R^{a-1}$ agrees with the
restriction of $P^a$ by Lemma~\ref{norm:naturality}.  The pushout
therefore supplies $f^a$, including the required boundary data.
Although the displayed pushout consists of $D$-algebras, it is also
the pushout of the same span of rational algebras: the common middle
algebra already carries the specified scalar maps.  Thus this
construction does not impose $D$-linearity on $P^a$.

Define the homogeneous cycle class and the next algebra by
\begin{equation}
\label{norm:next-cell}
 g_a=f^a(b_a^u)\in\pi_0(L_a)_{i_a},
 \qquad L_{a+1}=L_a//g_a.
\end{equation}
Use the attached null as $e'_a$, and put
$L_{a+1}^M=L_0^M\otimes_{L_0}L_{a+1}$.
The first equivalence in Equation~\eqref{eq:target-layer} carries
the reduction of $g_a$ to a class of
$h_\Omega(Y_{a+1})$ at weight $F(i_a)=w_{a+1}$.  Represent this
class by a strict cycle $b_{a+1}$ in $A[e_0,\ldots,e_a]$, choosing
a path matching its class with that reduction.  Attach the next
null $e_{a+1}$ and realize this finite-module algebra extension as
$Y_{a+2}$ by Proposition~\ref{found:koszul-realization}.

\begin{lemma}[Compatibility of the new attachment]
\label{norm:attachment}
The preceding choices extend Equation~\eqref{eq:target-layer} to
stage $a+1$, with its unit, initial-layer, and tower compatibilities.
Both new algebras are derived $t$-complete.
\end{lemma}

\begin{proof}
Algebra pushouts commute with one another, so reducing the first
attachment modulo $t$ gives
\[
 L_{a+1}//t\simeq (L_a//t)//\overline g_a.
\]
The old layer equivalence, the chosen path matching attaching
classes, and the actual null $e_{a+1}$ define an algebra map from
this pushout to $F^*h_\Omega(Y_{a+2})$.  We check this particular
map on its underlying module.

At an index $i\in I$, the source module is the cofiber of
multiplication by $\overline g_a$ from
$(L_a//t)_{i-i_a}$ to $(L_a//t)_i$.  Under the old equivalence these
terms become
\[
 h_\Omega(Y_{a+1})_{F(i)-w_{a+1}}
 \xrightarrow{\ b_{a+1}\ }
 h_\Omega(Y_{a+1})_{F(i)}.
\]
Their cofiber is $h_\Omega(Y_{a+2})_{F(i)}$.  The map on the new
cell is the chosen null $e_{a+1}$ multiplied by the base input.
Lemma~\ref{found:free-null} therefore identifies the constructed
map with this cofiber equivalence.  This argument uses only the
additivity of $F$ and a two-term cofiber at each index.  It does not
require $F$ to be injective or its pullback to preserve arbitrary
relative algebra tensors.

For the $M$-tower, base change gives
\[
 L_{a+1}^M//t\simeq
 (L_a^M//t)\otimes_{L_a//t}(L_{a+1}//t).
\]
Map its first input into $F^*h_\Omega(M\otimes Y_{a+2})$ by the
old $M$-comparison followed by the actual $Y$-attachment.  Map its
second input there by the new sphere comparison followed by the
$M$-unit.  Their agreement on $L_a//t$ is the previous unit
compatibility together with the same attaching-class path.
The resulting pushout map is again the underlying cofiber
comparison: the new null is the image of $e_{a+1}$ under the
$M$-unit, and the finite-module tensor comparison identifies the
target with the cofiber of multiplication by this image of
$b_{a+1}$.  This proves the second equivalence and its compatibility
with the first.  The construction uses the canonical maps from the
old algebras, so it also preserves the initial-layer and tower maps.

The new attachment is a finite perfect-module extension of its
base.  Completeness follows from Lemma~\ref{layer:completion}, for
both the sphere tower and its $M$-base change.
\end{proof}

Figure~\ref{fig:towers} displays the induction just proved. The
auxiliary tower is one attachment ahead after reduction: its $a$th
stage records the ordinary object $Y_{a+1}$. The comparison is
compatible with the specified nulls, which is essential for the
next cycle to belong to that ordinary object.
\begin{figure}[htbp]
\centering
\[
\begin{array}{ccc}
 L_a & \xrightarrow{\ \text{adjoin }e'_a\ } & L_{a+1} \\[5pt]
 \big\downarrow\mathrlap{\scriptstyle\,//t} &&
 \big\downarrow\mathrlap{\scriptstyle\,//t} \\[5pt]
 F^*h_\Omega(Y_{a+1}) & \longrightarrow & F^*h_\Omega(Y_{a+2})
\end{array}
\]
\caption{Reduction of the auxiliary tower constructs the next ordinary
attaching class. The upper arrow nulls $g_a$; under the vertical layer
equivalences its reduction is $b_{a+1}$, and the lower arrow adjoins
the specified null $e_{a+1}$. The same square holds after the
compatible base change to complex cobordism.}
\label{fig:towers}
\end{figure}


\begin{proposition}[The compatible towers]
\label{norm:construction}
There exist cycles $b_0,\ldots,b_n$ and their realizations
$Y_0,\ldots,Y_{n+1}$, together with the algebras $L_a,L_a^M$ for
$0\leq a\leq n$, satisfying Equation~\eqref{eq:target-layer} and
its sphere-to-$M$, initial-layer, tower-map, and ordinary
$\delta$-scalar compatibilities. Each $b_j$ has chain degree zero
and weight $w_j$. For each $0\leq a<n$
there are rational algebra maps $P^a$ and $f^a$ as constructed above,
and $g_a=f^a(b_a^u)$ has chain degree zero and weight $i_a$.
It defines the next tower attachment and reduces modulo $t$ to
$b_{a+1}$ under Equation~\eqref{eq:target-layer}.
\end{proposition}

\begin{proof}
The initial layer was established at the beginning of the section.
Given stage $a<n$, Propositions~\ref{universal:primewise-models} and
\ref{conn:k-test} supply the universal models and
Equation~\eqref{eq:k-test}.  Lemmas~\ref{norm:target-locality} and
\ref{norm:normalization} define $P^a$; Lemma~\ref{norm:naturality}
and the older attached nulls extend it to $f^a$.
Equation~\eqref{norm:next-cell} defines $g_a$, and
Lemma~\ref{norm:attachment} constructs the next cycle and preserves
the invariant.  There are exactly $n$ such steps.
\end{proof}

\subsection{The comparison used for cobordism detection}

The construction has produced the next cycle before identifying
its Hurewicz image.  We finish this section with the comparison that
will perform that identification.  It relates the normalized
operation on the universal cycle to the original power on $M$.

\begin{proposition}[Recovery on the preceding Koszul object]
\label{norm:recovery}
For $0\leq a<n$ there is a map of $\Lambda$-graded rational algebras
\begin{equation}
\label{eq:recovery}
 h_\Lambda(M\otimes Y_a)
 \simeq h_\Lambda(M)\otimes_{R^a}K_-^a
 \longrightarrow \bigl(L_a^M[t^{-1}]\bigr)^\flat.
\end{equation}
On $h_\Lambda(M)$ it is the unnormalized power $P_M$ followed by
the map to the completed tower and inversion.  The $M$-Hurewicz
image of $b_a$ maps to $t g_a^M$, where $g_a^M$ is the image of
$g_a$ in $L_a^M$.
\end{proposition}

\begin{proof}
The actual evaluation $R^a\to A|_\Lambda$ sends the universal
coefficients to those of the chosen cycles.  Consequently
\[
 \bigl(A[e_0,\ldots,e_{a-1}]\bigr)|_\Lambda
 \simeq A|_\Lambda\otimes_{R^a}K_-^a,
\]
with each universal null sent to its actual counterpart.  Tensoring
with $h_\Lambda(M)$ over $A|_\Lambda$ gives the equivalence on the
left of Equation~\eqref{eq:recovery}.  Its identification with
$h_\Lambda(M\otimes Y_a)$ follows from the finite-module tensor
comparison of Proposition~\ref{found:finite-cell-morita}.  Only the finite
Koszul module is used here, so the same comparison applies after
restriction from $\Omega$ to $\Lambda$.

Define a map on the first pushout input $h_\Lambda(M)$ using
$P_M$ and the maps
$B(M)[t^{-1}]\to L_0^M[t^{-1}]\to L_a^M[t^{-1}]$.
On the other input use
\[
 K_-^a\xrightarrow{f^a}L_a^\#\longrightarrow(L_a^M)^\#
        \xrightarrow{s_{L_a^M}}
                    (L_a^M[t^{-1}])^\flat.
\]
These maps agree on $R^a$.  Indeed, restriction of the second map
to $R^a$ recovers the evaluated, unnormalized power by
Equation~\eqref{norm:raw-recovery}.  Naturality of the primewise
power for $R_p^a\to S_{(p)}\to\MU_{(p)}$ identifies this with
$P_M$ applied to the evaluated coefficients.  The chosen path
between spectral evaluation and the actual coefficient evaluation
is part of this comparison.  It gives agreement as algebra maps,
including the scalar action and the boundary data.

The universal property of the pushout now gives the arrow in
Equation~\eqref{eq:recovery}.  Although its source was described
using $D$-algebras, the same span is a pushout of rational algebras:
the common algebra $R^a$ supplies the scalar identifications.  No
additional $D$-linearity condition on the map into the powered
target is required.

Finally, the Hurewicz image of $b_a$ corresponds in that pushout to
the universal cycle $b_a^u\in K_-^a$.  Its image under $f^a$ is
$g_a$ by definition.  Since $\nu(w_a)=1$, the map
$s_{L_a^M}$ multiplies this class by $t$, changing its index from
$(w_a,0,-1)$ to $(w_a,0,0)$.  This proves the asserted value
$t g_a^M$.
\end{proof}
